% Section 3.4, from lectures/L03/appendix/c_exercises.md.

\section{Exercises}
\label{c3:sec:exercises}\label{c3:app:c}

\begin{aside}
\textbf{How to work these.} Exercise~\ref{c3:ex:1} reinforces \secref{c3:sec:sync} and
Exercise~\ref{c3:ex:2} \secref{c3:sec:uart_rx}. Design each module from the specification in its
own section. The FIFO, the two instantiations that put the receive path into \code{uart_top}, and
the overrun question belong to \chapref{4} (\secref{c4:sec:exercises}), not to this chapter.

\textbf{How to check your work.} Check each module with its provided testbench, run from
\file{hw/}; the preface's \emph{Setting up} has the full flow.

The answers to every part that is not code are in \secref{sol:c3}. Try each part before you read its
answer.
\end{aside}

% ------------------------------------------------------------------------------------------
\exercise{\code{sync}}{VHDL}
\label{c3:ex:1}

\xpart{a} Write your own \code{sync.vhd} from the specification in \secref{c3:sec:sync} and run its
testbench:

\begin{shell}
cd hw
ghdl -a --std=93 sync.vhd sync_tb.vhd
ghdl -e --std=93 sync_tb
ghdl -r --std=93 sync_tb --assert-level=error
\end{shell}

\noindent Note that no package appears in that first line: \code{sync} reads none, which is part of
what makes it reusable.

\xpart{b} Delete one of the two flip-flops, leaving a single register, and re-run. Which check
fails, and what does its message tell you? Put it back. Why would this bug pass every \emph{other}
simulation in the course if this one testbench did not exist?

\xpart{c} \code{sync} clears on \code{reset_s2_n}, and the testbench asserts it \emph{between} two
clock edges and requires the output to clear before the next one. Change your reset to a
synchronous one, tested inside \code{if rising_edge(clock)}, and re-run. Which check fails, and what
is the message? Put it back, then say why the assertion has to be asynchronous while the
\emph{release} does not need to be, given where \code{reset_s2_n} comes from.

\xpart{d} \code{sync} takes \code{reset_s2_n}, the already-synchronized reset, as an input. Explain
why that one fact makes it impossible for \code{sync} to be the module that produces
\code{reset_s2_n}, and why \code{reset_sync} therefore exists as a separate provided block even
though both are two flops in series.

\xpart{e} \code{uart_top} uses \code{sync} at \code{COUNT = 1}. Two bits of a \code{COUNT = 8}
instance change on the same edge; explain why \code{sync_out} may show one bit's new value and the
other's old one for a clock, why that is not a bug in \code{sync}, and why it would matter if those
eight bits were a counter read across clock domains rather than eight unrelated pins.

% ------------------------------------------------------------------------------------------
\exercise{\code{uart_rx}}{VHDL}
\label{c3:ex:2}

\xpart{a} Write your own \code{uart_rx.vhd} from the specification in \secref{c3:sec:uart_rx} and
run its testbench:

\begin{shell}
cd hw
ghdl -a --std=93 uart_def.vhd uart_rx.vhd uart_rx_tb.vhd
ghdl -e --std=93 uart_rx_tb
ghdl -r --std=93 uart_rx_tb --assert-level=error
\end{shell}

\noindent \code{uart_def.vhd} comes first because \code{uart_rx_tb} reads the package, for the
\code{to_hex} in its failure messages. \code{sync.vhd} is not on that line at all: \code{uart_rx}
reads \code{rx_s2}, the already-synchronized line, and instantiates nothing. The crossing belongs
one level up, in \code{uart_top}, which is where Exercise~\ref{c4:ex:2} puts it.

\xpart{b} The receiver samples each bit at tick 8 of 16. Suppose the transmitter's baud rate is a
few percent faster than the receiver's, so the receiver's ticks are slightly long. Which end of the
byte drifts furthest from the bit centre, the first data bit or the stop bit, and why? Roughly how
much faster could the transmitter run before the stop-bit sample landed outside its bit?

\xpart{c} The \code{STATE_START} state re-checks \code{rx_s2} at tick 8 before committing to a
frame. Remove that check (go straight from the falling edge to \code{STATE_DATA}) and describe what
a single-tick glitch on an idle line would now produce. Which real-world condition makes this more
than a theoretical worry?

\xpart{d} Trace the store \code{frame(bit_idx) <= rx_s2} for the byte \code{0x53}
(\code{0101_0011}) sent least significant first. Write the eight received bits in the order they
arrive, name the index each one lands in, and show that after eight bits \code{frame} holds
\code{0x53}. Then: what would \code{data_out} read if \code{bit_idx} counted \textbf{down} from 7
instead of up from 0, and what is the relationship between that byte and the one that was sent?
